\documentclass{article}
\usepackage{algorithm, algorithmic, amsfonts, amsmath}

\begin{document}
\begin{algorithm}[!h]
\renewcommand{\algorithmicrequire}{\textbf{Input:}}
\renewcommand{\algorithmicensure}{\textbf{Output:}}
	\caption{Presence-only EM}
	\begin{algorithmic}[1]
	    \REQUIRE
	    Covariates $\mathbf{X}$; observation indicator $z_{i}$ ($1$ = recorded presence,
	    $0$ = background site); population prevalence $\pi=P(y=1)$;
	    tolerance $\epsilon$
	    \ENSURE
	    Fitted model and the corrected linear predictor $\hat{\eta}$
	    \STATE \textbf{Setting:} true presence $y$ is observed only when $z=1$; background sites
	    have unknown $y$, so $y$ is treated as missing data % absence is very hard to confirm in ecology, so only presence and background points exist
	    \STATE $\hat{y}_{i}^{(0)}\leftarrow 1$ for $z_{i}=1$;\quad
	    $\hat{y}_{i}^{(0)}\leftarrow \pi$ for $z_{i}=0$ % before any model is fitted, the best guess at a background point is the unconditional prior
	    \STATE $\mathrm{adj}\leftarrow \log\!\left(\dfrac{n_{p}+\pi n_{u}}{\pi n_{u}}\right)$,
	    with $n_{p}=\#\{z_{i}=1\}$, $n_{u}=\#\{z_{i}=0\}$
	    \FOR{$k=1,2,\dots$ until convergence}
	    \STATE \textbf{M-step:} fit a standard logistic model (or GAM, boosted trees) of
	    $\hat{\mathbf{y}}^{(k-1)}$ on $\mathbf{X}$, giving $\hat{\eta}^{\ast(k)}$ % the response is a real number in [0,1], so a non-integer label needs the weighted-observation treatment
	    \STATE \quad $\hat{\eta}^{(k)}\leftarrow \hat{\eta}^{\ast(k)}-\mathrm{adj}$
	    \hfill \emph{case--control intercept correction} % background points are sampled at a different rate from the real landscape, so the intercept must be corrected
	    \STATE \textbf{E-step:}
	    $\hat{y}_{i}^{(k)}\leftarrow \dfrac{e^{\hat{\eta}^{(k)}(\mathbf{x}_{i})}}{1+e^{\hat{\eta}^{(k)}(\mathbf{x}_{i})}}$
	    for $z_{i}=0$;\quad $\hat{y}_{i}^{(k)}\leftarrow 1$ for $z_{i}=1$
	    \STATE Stop when $\max_{i}\left|\hat{y}_{i}^{(k)}-\hat{y}_{i}^{(k-1)}\right|<\epsilon$
	    \ENDFOR
		\RETURN the fitted model and $\hat{\eta}$
	\end{algorithmic}
\end{algorithm}
% Ward, Hastie, Barry, Elith, Leathwick, Biometrics 2009,
% "Presence-Only Data and the EM Algorithm".
% The application is species distribution modelling: only locations where the species was observed, plus background
% points drawn at random from the whole landscape, whose true occupancy is unknown. That is the PU structure.
% Why branch C: the population prevalence pi enters the model as a known quantity, both initializing the response at
% background points and correcting the sampling bias through the intercept term log((n_p + pi n_u)/(pi n_u)).
%
% Three implementation notes given in the original:
% (1) the y_hat produced by the E step is a real number in [0,1] at background points, and where a fitting routine
%     rejects non-integer responses, weighted observations or an equivalent split-data device can be used;
% (2) for stepwise algorithms such as boosted trees, which need hundreds of steps, the E step can be interleaved into
%     tree construction (re-estimating y_hat every ten trees, say), which speeds convergence by about an order of magnitude;
% (3) identifiability. Unless eta(x) is strongly constrained, for instance to be strictly linear, pi is not
%     identifiable from presence-only data. It cannot be estimated inside the algorithm and must be supplied from
%     external data or expert knowledge, with a sensitivity analysis over pi.
% Point (3) matches the identifiability thread running through the survey: the prior under SCAR and the propensity
% under SAR are both quantities the data cannot settle and that an outside assumption must supply.
\end{document}
